\section*{Week 3 assignment}


\question{1}{}
Let $O(n)=\{A\in M_n(\mathbb{R}) \mid A^TA=I\}$ be the orthogonal group, where $A^T$ is the transpose of $A$. Consider the map
\begin{align}
    \varphi: M_n(\mathbb{R})\to \mathrm{Sym}(n),\quad A\mapsto A^TA,
\end{align}
where $\mathrm{Sym}(n)$ is the set of symmetric $n\times n$ matrices (i.e., $B^T=B$).
\begin{enumerate}
    \item Show that $\mathrm{Sym}(n)$ is a manifold. Compute its dimension.
    \item Compute the derivative of $\varphi$ and show that $\varphi$ is a submersion.
    \item Prove that $O(n)$ is a submanifold of $M_n(\mathbb{R})$. What is the dimension of $O(n)$?
    \item Prove that $O(n)$ is compact.
\end{enumerate}
\answer{}

\begin{namedthm}{Definition 7.1}[Submersion, Ko Honda's notes]
Let $U\subseteq\mathbb{R}^m$ and $V\subseteq\mathbb{R}^n$ be open sets. A smooth map $f:U\to V$ is a \textbf{submersion} if $df(x)$ is surjective for all $x\in U$. More generally, a smooth map $f:M\to N$ between manifolds is a submersion if all of its local representatives $\psi_\beta\circ f\circ\phi_\alpha^{-1}$ are submersions.
\end{namedthm}

\begin{namedthm}{Corollary 7.4}[Regular value theorem, Ko Honda's notes]
Let $f: M\to N$ be a smooth map, where $\dim M = m$ and $\dim N = n$. A point $y\in N$ is called a \textbf{regular value} of $f$ if $df(x)$ is surjective for every $x\in f^{-1}(y)$. If $y$ is a regular value of $f$, then $f^{-1}(y)$ is a submanifold of $M$ of dimension $m-n$.
\end{namedthm}




\begin{enumerate}
    \item [1.]
\end{enumerate}
Consider the map $\psi:\mathbb{R}^{n(n+1)/2}\to \mathrm{Sym}(n)$ defined by 
\begin{align}
    \psi(x_1,\cdots,x_{n(n+1)/2})=\begin{pmatrix}
        x_1 & x_2 & \cdots & x_n\\
        x_2 & x_{n+1} & \cdots & x_{2n-1}\\
        \vdots & \vdots & \ddots & \vdots\\
        x_n & x_{2n-1} & \cdots & x_{n(n+1)/2}
    \end{pmatrix}.
\end{align}
This map is a bijection, and its inverse is given by
\begin{align}
    \psi^{-1}\left(\begin{pmatrix}
        a_{11} & a_{12} & \cdots & a_{1n}\\
        a_{12} & a_{22} & \cdots & a_{2n}\\
        \vdots & \vdots & \ddots & \vdots\\
        a_{1n} & a_{2n} & \cdots & a_{nn}
    \end{pmatrix}\right)=(a_{11},a_{12},\cdots,a_{1n},a_{22},a_{23},\cdots,a_{2n},\cdots,a_{nn}).
\end{align}
We equip $\mathrm{Sym}(n)$ with the subspace topology from $M_n(\mathbb{R})\cong\mathbb{R}^{n^2}$. Since $\psi$ is linear, it is continuous, and $\psi^{-1}$ is the restriction of a linear projection $M_n(\mathbb{R})\to\mathbb{R}^{n(n+1)/2}$, hence also continuous. So $\psi^{-1}$ is a homeomorphism.

Hence $\{(\mathrm{Sym}(n),\psi^{-1})\}$ is an atlas consisting of a single chart. The only transition map is $\psi^{-1}\circ\psi=\mathrm{id}$, which is smooth.

Since $\mathrm{Sym}(n)$ is a subspace of $\mathbb{R}^{n^2}$, which is Hausdorff and second countable, $\mathrm{Sym}(n)$ is also Hausdorff and second countable.

Therefore, $\mathrm{Sym}(n)$ is a manifold of dimension $n(n+1)/2$.

\begin{enumerate}
    \item [2.]
\end{enumerate}
Consider the limit 
\begin{align}
    &\lim_{t\to 0}\frac{\varphi(A+tB)-\varphi(A)}{t}\\
    =&\lim_{t\to 0}\frac{(A+tB)^T(A+tB)-A^TA}{t}\\
    =&\lim_{t\to 0}\frac{A^TA+t(B^TA+A^TB)+t^2B^TB-A^TA}{t}=B^TA+A^TB.
\end{align}
It is easy to see that $B^T A+A^TB$ is symmetric. Then the derivative of $\varphi$ at $A$ is given by
\begin{align}
    d\varphi(A):M_n(\mathbb{R})\to \mathrm{Sym}(n),\quad B\mapsto B^TA+A^TB.
\end{align}
We show that $d\varphi(A)$ is surjective for every $A\in O(n)$. Consider $A\in O(n)$, then $A^TA=I$. Let $C\in\mathrm{Sym}(n)$ be arbitrary. We want to find $B\in M_n(\mathbb{R})$ such that $B^TA+A^TB=C$. Choose $B=\frac{1}{2}AC$, then we have
\begin{align}
    d\varphi(A)(B)=B^TA+A^TB=\frac{1}{2}(AC)^TA+\frac{1}{2}A^T(AC)=\frac{1}{2}C^TA^TA+\frac{1}{2}A^TAC=\frac{1}{2}C^T+\frac{1}{2}C=C.
\end{align}
Hence $d\varphi(A)$ is surjective for every $A\in O(n)=\varphi^{-1}(I)$, i.e., $I$ is a regular value of $\varphi$. 

To be more clear, $\varphi$ is a submersion at $A\in O(n)$, but not a submersion at $0\in M_n(\mathbb{R})$, since $d\varphi(0)=0$.

\begin{enumerate}
    \item [3.]
\end{enumerate}
By definition, $O(n)=\varphi^{-1}(I)$. Since $I$ is a regular value of $\varphi$, by Corollary 7.4, $O(n)$ is a submanifold of $M_n(\mathbb{R})$ of dimension
\begin{align}
    \dim O(n)=\dim M_n(\mathbb{R})-\dim \mathrm{Sym}(n)=n^2-\frac{n(n+1)}{2}=\frac{n(n-1)}{2}.
\end{align} 

\begin{enumerate}
    \item [4.]
\end{enumerate}
To show that $O(n)$ is compact, we need to show that it is closed and bounded by Heine-Borel theorem. Since $O(n)\subset M_n(\mathbb{R})\cong\mathbb{R}^{n^2}$, it suffices to show that $O(n)$ is closed and bounded in $\mathbb{R}^{n^2}$.

Since $\{I\}$ is closed in $\mathrm{Sym}(n)$ and $\varphi$ is continuous, $O(n)=\varphi^{-1}(\{I\})$ is closed.

To show that $O(n)$ is bounded, we need to show that there exists $M>0$ such that $\|A\|\le M$ for all $A\in O(n)$. Consider the Frobenius norm $\|A\|_F=\sqrt{\sum_{i,j}a_{ij}^2}=\sqrt{\mathrm{Tr}(A^TA)}$. Since $A\in O(n)$, we have $A^TA=I$, so $\|A\|_F=\sqrt{\mathrm{Tr}(I)}=\sqrt{n}$. Therefore, $O(n)$ is bounded.

Hence, $O(n)$ is compact.\qed

\clearpage
\question{2}{}
Let $F(x_1,\cdots,x_n)$ be a homogeneous function of degree $d$ in $n$ real variables, i.e.,
\begin{align}
    F(tx_1,\cdots,tx_n)=t^d\cdot F(x_1,\cdots,x_n).
\end{align}
\begin{enumerate}
    \item Prove Euler's identity: $\displaystyle\sum_{i=1}^{n}x_i\frac{\partial F}{\partial x_i}=d\cdot F$.
    \item Prove that the set $F^{-1}(a)$, for $a\neq 0$, is a submanifold of $\mathbb{R}^n$.
\end{enumerate}
\answer{}



\begin{enumerate}
    \item [1.]
\end{enumerate}
Consider the function $g(t)=F(tx_1,\cdots,tx_n)$. Then $g(t)=t^d\cdot F(x_1,\cdots,x_n)$. Taking the derivative with respect to $t$, we have
\begin{align}
    g'(t)=\sum_{i=1}^{n}\frac{\partial F}{\partial x_i}(tx_1,\cdots,tx_n)\cdot x_i=dt^{d-1}\cdot F(x_1,\cdots,x_n).
\end{align}
Setting $t=1$, we get
\begin{align}
    \sum_{i=1}^{n}x_i\frac{\partial F}{\partial x_i}(x_1,\cdots,x_n)=d\cdot F(x_1,\cdots,x_n).
\end{align} 

\begin{enumerate}
    \item [2.]
\end{enumerate}
Assume $F$ is smooth and $d>0$. To show that $F^{-1}(a)$ is a submanifold of $\mathbb{R}^n$, we show that $a$ is a regular value of $F$, i.e., $dF(x)$ is surjective for every $x\in F^{-1}(a)$.

Let $x\in F^{-1}(a)$, so $F(x)=a$. For $v\in\mathbb{R}^n$,
\begin{align}
    dF(x)(v)=\lim_{t\to 0}\frac{F(x+tv)-F(x)}{t}=\sum_{i=1}^{n}\frac{\partial F}{\partial x_i}(x)\cdot v_i.
\end{align}
Let $w\in\mathbb{R}$ be arbitrary. We want to find $v\in\mathbb{R}^n$ such that $dF(x)(v)=w$. First, taking $v=x$, by Euler's identity we have
\begin{align}
    dF(x)(x)=\sum_{i=1}^{n}x_i\frac{\partial F}{\partial x_i}(x)=d\cdot F(x)=d\cdot a\neq 0.
\end{align}
Hence, for $v=\frac{w}{d\cdot a}\,x$, by linearity of $dF(x)$,
\begin{align}
    dF(x)(v)=\frac{w}{d\cdot a}\,dF(x)(x)=\frac{w}{d\cdot a}\cdot d\cdot a=w.
\end{align}
Hence $dF(x)$ is surjective for every $x\in F^{-1}(a)$, i.e., $a$ is a regular value of $F$. By Corollary 7.4, $F^{-1}(a)$ is a submanifold of $\mathbb{R}^n$ of dimension $\dim\mathbb{R}^n-\dim\mathbb{R}=n-1$.\qed

\textbf{Remark.} The assumption $d\neq 0$ is needed. If $d=0$, then $F(tx)=F(x)$ for all $t$; letting $t\to 0$ and using the continuity of $F$ gives $F(x)=F(0)$, so $F$ is constant. Then $F^{-1}(a)$ is either empty or all of $\mathbb{R}^n$, which is not of dimension $n-1$.